\vspace{0.5em}\noindent\textbf{Kiến trúc}\par

\vspace{0.5em}\noindent\textbf{Mục tiêu}\par

Chương tài liệu này tường giải cấu trúc kiến trúc AWS sau cùng cũng như
các luồng thực thi (runtime flows) trọng tâm của hệ thống Internship
Application Tracker.

\vspace{0.5em}\noindent\textbf{Phạm vi}\par

Kiến trúc này được phác thảo dựa trên minh chứng production sau cùng,
tài liệu thiết kế ghi chú của dự án và kho lưu trữ mã nguồn hiện tại của
ứng dụng. Tôi lấy bằng chứng vận hành khi chạy thực tế (runtime
evidence) làm cội nguồn chân lý cho các tình huống mà nó có điểm xô lệch
so với những manifest cũ. Điểm đính chính kiến trúc đáng chú ý nhất hiện
nay là việc frontend được lưu giữ và phát tán qua cống bộ đôi S3 và
CloudFront, loại bỏ yếu tố chạy trong EKS.

\vspace{0.5em}\noindent\textbf{Bối cảnh kiến
trúc}\par

{
\begingroup
\small
\setlength{\tabcolsep}{3pt}
\renewcommand{\arraystretch}{1.15}
\begin{longtable}{@{}
  >{\raggedright\arraybackslash}p{(\linewidth - 2\tabcolsep) * \real{0.5000}}
  >{\raggedright\arraybackslash}p{(\linewidth - 2\tabcolsep) * \real{0.5000}}@{}}
\toprule\noalign{}
\begin{minipage}[b]{\linewidth}\raggedright
Lớp (Layer)
\end{minipage} & \begin{minipage}[b]{\linewidth}\raggedright
Các thành phần chính (Main components)
\end{minipage} \\
\midrule\noalign{}
\endhead
\bottomrule\noalign{}
\endlastfoot
Ngoại vi (Edge) & Tài nguyên CloudFront distribution
\texttt{EQIGYNECXDYL8} \\
Định tuyến công cộng (Public routing) & Cống điều phối lưu lượng
Application Load Balancer
\texttt{k8s-\allowbreak{}internshippublic-\allowbreak{}48101b50ad-\allowbreak{}85486086.\allowbreak{}ap-\allowbreak{}southeast-\allowbreak{}1.\allowbreak{}elb.\allowbreak{}amazonaws.\allowbreak{}com} \\
Cụm Kubernetes & Cụm máy EKS \texttt{internship-\allowbreak{}prod}, không gian tên
\texttt{internship} \\
Khối workload ứng dụng & \texttt{backend}, \texttt{chat-\allowbreak{}service},
\texttt{backend-\allowbreak{}outbox-\allowbreak{}dispatcher}, \texttt{backend-\allowbreak{}processing-\allowbreak{}worker},
\texttt{ai-\allowbreak{}service} \\
Lớp dữ liệu (Data) & Cỗ máy RDS PostgreSQL, DynamoDB, ElastiCache Redis,
kho object storage S3 \\
Lớp bất đồng bộ và sự kiện (Async and event) & Hàng đợi công việc trong
PostgreSQL, bảng PostgreSQL outbox, hàng đợi SQS, hàm Lambda, dịch vụ
phát thư SES \\
Trí tuệ nhân tạo (AI) & Pod trung gian \texttt{ai-\allowbreak{}service} adapter kết
cọc cùng trạm suy đoán SageMaker endpoint
\texttt{internship-\allowbreak{}qwen3-\allowbreak{}4b} \\
Dây chuyền CI/CD & GitHub Actions, ủy quyền AWS OIDC, kho ảnh ECR, công
cụ kubectl, bộ lệnh AWS CLI \\
\end{longtable}
\endgroup
}

\vspace{0.5em}\noindent\textbf{Sơ đồ 1: Tổng thể kiến trúc
AWS}\par

\textbf{Mô tả sơ bộ.} Người dùng truy cập qua CloudFront; CloudFront
phục vụ frontend S3 riêng tư hoặc chuyển yêu cầu động đến ALB và EKS.
Các workload EKS kết nối RDS, DynamoDB, Redis, S3, SageMaker và tuyến
thông báo SQS/Lambda; DLQ cùng bảng DynamoDB dedupe hỗ trợ xử lý sự kiện
tin cậy.

\vspace{0.5em}\noindent\textbf{Sơ đồ 2: Luồng định tuyến yêu cầu trên
CloudFront}\par

\textbf{Mô tả sơ bộ.} CloudFront kiểm tra đường dẫn của từng yêu cầu
HTTPS. Luồng mặc định phục vụ frontend từ S3; \texttt{/\allowbreak{}api/\allowbreak{}*} chuyển đến
backend service; \texttt{/\allowbreak{}chat/\allowbreak{}*} và \texttt{/\allowbreak{}socket.\allowbreak{}io/\allowbreak{}*} chuyển đến
chat service thông qua ALB.

Kho lưu trữ bucket cho frontend duy trì an tọa trong tư thế khép kín
(private). CDN CloudFront ôm gánh trách nhiệm mã hóa HTTPS và điều luật
đường dự phòng cho giao diện trang đơn (SPA fallback). Luồng ALB Ingress
thực thi cắt gõ xóa phần tiền tố \texttt{/\allowbreak{}api} cùng \texttt{/\allowbreak{}chat}
(rewrites) trước thời khắc luân chuyển thẳng về cho 2 dịch vụ backend và
chat.

\vspace{0.5em}\noindent\textbf{Sơ đồ 3: Biểu đồ tuần tự luồng tải CV và tác vụ xử lý
AI}\par

\textbf{Mô tả sơ bộ.} Ứng viên nộp CV qua frontend; backend lưu tham
chiếu tài liệu cùng các bản ghi ứng tuyển, idempotency, outbox và
processing. Worker nhận job đang chờ, gọi AI adapter và SageMaker, lưu
kết quả chuẩn hóa vào PostgreSQL, sau đó frontend truy vấn trạng thái để
nhận kết quả.

\vspace{0.5em}\noindent\textbf{Sơ đồ 4: Biểu đồ tuần tự luồng tin nhắn chat thời gian thực
(Realtime
chat)}\par

\textbf{Mô tả sơ bộ.} Lưu lượng Socket.IO đi qua CloudFront và ALB đến
một pod chat-service. Pod nhận tin nhắn sẽ lưu vào DynamoDB, phát sự
kiện phòng qua Redis để các pod khác thông báo cho người dùng đang theo
dõi, rồi mới gửi xác nhận cho người gửi.

Cơ sở dữ liệu DynamoDB là cõi lưu cất vững chắc vĩnh cửu cho tin nhắn
chat. Cụm ElastiCache Redis duy chỉ làm nhiệm vụ trạm chuyển nhĩ cẩu
pub/sub giao thoa giữa các pod và hoàn toàn không xô lấn thay thế thế
trận lưu dữ liệu bền lâu (persistent storage).

\vspace{0.5em}\noindent\textbf{Sơ đồ 5: Biểu đồ tuần tự luồng transactional outbox, SQS và
Lambda}\par

\textbf{Mô tả sơ bộ.} Backend ghi thay đổi nghiệp vụ và sự kiện outbox
trong một transaction PostgreSQL. Dispatcher gửi sự kiện vào SQS; Lambda
dùng conditional write trên DynamoDB để khử trùng lặp. Sự kiện lần đầu
được lưu trữ ở S3 và có thể gửi email qua SES, còn sự kiện trùng được
xem là đã xử lý.

Cơ quan vận chuyển tin SQS bám sát cam kết phân phối ít nhất một lần
(at-least-once). Việc xuất hiện các nhịp tin chao bị nhận ban lặp lại là
điều được tiên định trong thiết kế, do vậy khả năng xử lý bất biến
(idempotency) của Lambda là nguyên tắc sống còn. Kịch bản test khói
(smoke-test) ghi trát thành quả đắc thắng đã cất mã định danh sự kiện là
\texttt{lambda-\allowbreak{}smoke-\allowbreak{}fixed-\allowbreak{}1785220478}, quy tập chôn an tọa dưới con
rãnh S3
\texttt{outbox-\allowbreak{}archive/\allowbreak{}2026/\allowbreak{}07/\allowbreak{}28/\allowbreak{}lambda-\allowbreak{}smoke-\allowbreak{}fixed-\allowbreak{}1785220478.\allowbreak{}json}.

\vspace{0.5em}\noindent\textbf{Sơ đồ 6: Luồng triển khai CI/CD khi thay đổi mã
nguồn}\par

\textbf{Mô tả sơ bộ.} Workflow GitHub Actions được kích hoạt thủ công để
kiểm tra repository và khôi phục compute khi cần, sau đó build image ECR
gắn thẻ SHA. Pipeline triển khai ứng dụng EKS và ALB, bảo đảm cấu hình
S3/CloudFront riêng tư, xuất bản frontend, tạo CloudFront invalidation
và tổng hợp kết quả triển khai.

Dây chuyền kịch bản workflow linh hoạt cung cẩu dải chế độ thi hành gồm
\texttt{validate}, \texttt{deploy}, \texttt{rollout},
\texttt{restore-\allowbreak{}compute}, \texttt{deploy-\allowbreak{}app}, \texttt{deploy-\allowbreak{}public},
\texttt{deploy-\allowbreak{}frontend}, và nhĩ lệnh cất còi trọn cõi \texttt{full}.
Trình GitHub Actions thông quan xác nhận thông tính ủy quyền AWS role
thông qua giao dịch OIDC; những chìa khóa bảo mật kiên bám AWS access
keys dài hạn tuyệt đối không có vị trí nào trong quy hoạch thiết lập an
sinh ban bố này.

\vspace{0.5em}\noindent\textbf{Mô tả các thành
phần}\par

{
\begingroup
\small
\setlength{\tabcolsep}{3pt}
\renewcommand{\arraystretch}{1.15}
\begin{longtable}{@{}
  >{\raggedright\arraybackslash}p{(\linewidth - 4\tabcolsep) * \real{0.3333}}
  >{\raggedright\arraybackslash}p{(\linewidth - 4\tabcolsep) * \real{0.3333}}
  >{\raggedright\arraybackslash}p{(\linewidth - 4\tabcolsep) * \real{0.3333}}@{}}
\toprule\noalign{}
\begin{minipage}[b]{\linewidth}\raggedright
Thành phần (Component)
\end{minipage} & \begin{minipage}[b]{\linewidth}\raggedright
Trách nhiệm (Responsibility)
\end{minipage} & \begin{minipage}[b]{\linewidth}\raggedright
Lệnh kiểm tra
\end{minipage} \\
\midrule\noalign{}
\endhead
\bottomrule\noalign{}
\endlastfoot
CloudFront & Cổng vào HTTPS công cộng cho người dùng đi đôi khả năng
điều phối phân tuyến &
\texttt{aws\ cloudfront\ get-\allowbreak{}distribution\ -\/\allowbreak{}-id\ EQIGYNECXDYL8} \\
S3 frontend bucket & Kho cất các tệp tài sản tĩnh của frontend &
\texttt{aws\ s3\ ls\ s3:/\allowbreak{}/\allowbreak{}internship-\allowbreak{}prod-\allowbreak{}frontend-\textless{}AWS\_\allowbreak{}ACCOUNT\_\allowbreak{}ID\textgreater{}/\allowbreak{}} \\
ALB & Luân chuyển điều trút nhịp truy vãng API/chat/socket rẽ theo hướng
vào các cụm service trên EKS &
\texttt{kubectl\ get\ ingress\ internship-\allowbreak{}public\ -n\ internship} \\
Backend & Cánh cổng nghiệp vụ FastAPI kinh doanh kiêm cõi điểm gọi sức
khỏe &
\texttt{kubectl\ rollout\ status\ deployment/\allowbreak{}backend\ -n\ internship} \\
Dịch vụ chat (Chat service) & Xử trí giao thông realtime Socket.IO kết
cọc cùng cõi chat REST API &
\texttt{kubectl\ rollout\ status\ deployment/\allowbreak{}chat-\allowbreak{}service\ -n\ internship} \\
Processing worker & Gác nhiệm vụ rèn giải mác hồ sơ và cọc AI theo cơ
cấu xử lý bất đồng bộ &
\texttt{kubectl\ get\ deployment\ backend-\allowbreak{}processing-\allowbreak{}worker\ -n\ internship} \\
Outbox dispatcher & Chiết gặt nhĩ tệp sự kiện từ cơ quan PostgreSQL để
bạt nạp sa cọc sang SQS &
\texttt{kubectl\ logs\ deployment/\allowbreak{}backend-\allowbreak{}outbox-\allowbreak{}dispatcher\ -n\ internship} \\
Dịch vụ AI (AI service) & Pod trung gian (SageMaker adapter) buông đường
vãng tuệ bền lâu cho đàn worker ra lời gọi &
\texttt{kubectl\ port-\allowbreak{}forward\ service/\allowbreak{}ai-\allowbreak{}service\ 8010:8010\ -n\ internship} \\
RDS PostgreSQL & Kho lưu giữ dữ liệu cho giao dịch và dải thông hàng đợi
&
\texttt{aws\ rds\ describe-\allowbreak{}db-\allowbreak{}instances\ -\/\allowbreak{}-db-\allowbreak{}instance-\allowbreak{}identifier\ internship-\allowbreak{}prod-\allowbreak{}postgres\ -\/\allowbreak{}-region\ ap-\allowbreak{}southeast-\allowbreak{}1} \\
DynamoDB & Nơi an trú các bảng chat cùng bàn lưu cất chứng tích khử lặp
cho Lambda &
\texttt{aws\ dynamodb\ describe-\allowbreak{}table\ -\/\allowbreak{}-table-\allowbreak{}name\ ChatMessages\ -\/\allowbreak{}-region\ ap-\allowbreak{}southeast-\allowbreak{}1} \\
Redis & Cơ quan chuyền tin phát chóp pub/sub mạn Socket.IO &
\texttt{aws\ elasticache\ describe-\allowbreak{}replication-\allowbreak{}groups\ -\/\allowbreak{}-replication-\allowbreak{}group-\allowbreak{}id\ internship-\allowbreak{}prod-\allowbreak{}redis\ -\/\allowbreak{}-region\ ap-\allowbreak{}southeast-\allowbreak{}1} \\
SQS & Con hào cống vận tải cho dòng chảy tin outbox &
\texttt{aws\ sqs\ get-\allowbreak{}queue-\allowbreak{}attributes\ -\/\allowbreak{}-queue-\allowbreak{}url\ \textless{}OUTBOX\_\allowbreak{}QUEUE\_\allowbreak{}URL\textgreater{}\ -\/\allowbreak{}-attribute-\allowbreak{}names\ All} \\
\end{longtable}
\endgroup
}

\vspace{0.5em}\noindent\textbf{Các tài nguyên công khai và riêng
tư}\par

{
\begingroup
\small
\setlength{\tabcolsep}{3pt}
\renewcommand{\arraystretch}{1.15}
\begin{longtable}{@{}
  >{\raggedright\arraybackslash}p{(\linewidth - 2\tabcolsep) * \real{0.5000}}
  >{\raggedright\arraybackslash}p{(\linewidth - 2\tabcolsep) * \real{0.5000}}@{}}
\toprule\noalign{}
\begin{minipage}[b]{\linewidth}\raggedright
Hướng ra Internet (Public-facing)
\end{minipage} & \begin{minipage}[b]{\linewidth}\raggedright
Riêng tư hoặc nội bộ (Private or internal)
\end{minipage} \\
\midrule\noalign{}
\endhead
\bottomrule\noalign{}
\endlastfoot
Tên miền công khai của CloudFront & Máy chủ CSDL RDS PostgreSQL \\
Tên miền DNS origin thuộc bãi ALB & Cụm máy nhân bản ElastiCache
Redis \\
Chuỗi các ranh định tuyến HTTPS của CloudFront & Hệ thống các bàn dữ
liệu bên DynamoDB \\
Các tài sản web tĩnh được trình duyệt tiếp thu qua CloudFront & Giao
thông luân chuyển thẳng từ pod tới service cõi EKS \\
Khâu ném phát đi văn thư mail qua cống SES & Hàng đợi sự kiện chủ lực
SQS đi đôi cọ bàn lỗi DLQ \\
& Cõi hòm kho cất S3 bị đóng tịt, khước bẻ trần truy cập trực diện lộ
thiên từ public \\
\end{longtable}
\endgroup
}

\vspace{0.5em}\noindent\textbf{Ranh giới bảo mật (Security
boundaries)}\par

\begin{itemize}
\tightlist
\item
  Tiến trình GitHub Actions nương tuệ nhờ quyền OIDC để ra chiêu hoán
  quyền thâu cọc vào \texttt{internship-\allowbreak{}github-\allowbreak{}deploy}.
\item
  Cụm các workload trên EKS vận dụng cấu hình ServiceAccount
  \texttt{internship-\allowbreak{}app} mang kèm theo bản gõ trói vai trò IRSA role
  mapping.
\item
  Muôn thông ký mật mã bí ẩn thọ an tiêm truyền thẳng cẩu từ hòm bí mật
  của GitHub và cọc Kubernetes Secret mang mác
  \texttt{internship-\allowbreak{}secrets}.
\item
  Kho chứa tệp tĩnh cho frontend trên S3 kỵ kiêng giêm giấu vùng kín
  riêng tư, duy nhĩ thọ ban visa mở cống qua duy nhất rào CDN
  CloudFront.
\item
  Các pod của backend và dịch vụ chat thọ nhĩ bọc rẽ khe hở duy trì cọc
  cõi giao thiệp ứng với ranh luồn cổng service nhờ vào luật chia đường
  của ALB.
\item
  Giao thông gọi dịch vụ AI trốn tiệm êm gọn trong ranh giới nội cluster
  đi từ pod worker tiến sang con pod \texttt{ai-\allowbreak{}service}; cõi trạm suy
  bạt SagMaker endpoint thọ khấn giục bục từ chính thợ bộ đệm adapter
  này.
\item
  Nghiệp vụ giao đãi giữa SQS/Lambda bấu rào cản tính bất biến, thiêng
  cọc bảo lưu qua ranh lệnh chốt ghi theo điều kiện trên cỗi DynamoDB.
\end{itemize}

\vspace{0.5em}\noindent\textbf{Kết quả mong đợi}\par

Cấu trúc kiến trúc trọn vẹn sau cùng thực hiện bóc rẽ trành minh tường
chuỗi phát tán web tĩnh, định tuyến truy thu từ công cộng, các dịch vụ
gác chạy mải khôn dừng, cõi cất trữ giao dịch có ACID, trạm đồng bộ tín
liệu realtime, rỗng bão chuyển tải sự kiện bất đồng bộ, cọc cống tin tức
chuế ngắn rêu cảnh báo, và cơ quan cất trát suy đoán AI. Trọn bộ từng
thực thể sắm lấy một cương vị trách nhiệm tường minh cùng khả năng khảo
rà, minh chứng hoạt vãng độc lập.

\vspace{0.5em}\noindent\textbf{Các lỗi thường
gặp}\par

{
\begingroup
\small
\setlength{\tabcolsep}{3pt}
\renewcommand{\arraystretch}{1.15}
\begin{longtable}{@{}
  >{\raggedright\arraybackslash}p{(\linewidth - 4\tabcolsep) * \real{0.3333}}
  >{\raggedright\arraybackslash}p{(\linewidth - 4\tabcolsep) * \real{0.3333}}
  >{\raggedright\arraybackslash}p{(\linewidth - 4\tabcolsep) * \real{0.3333}}@{}}
\toprule\noalign{}
\begin{minipage}[b]{\linewidth}\raggedright
Lỗi (Error)
\end{minipage} & \begin{minipage}[b]{\linewidth}\raggedright
Nguyên nhân
\end{minipage} & \begin{minipage}[b]{\linewidth}\raggedright
Hướng khắc phục
\end{minipage} \\
\midrule\noalign{}
\endhead
\bottomrule\noalign{}
\endlastfoot
Phác trát miêu tả frontend thọ chay nghênh gục bên giữa lòng EKS & Rơi
nhầm vào văn bản phác tài liệu kiến trúc thảm xa xưa & Nhẫn tâm bấu
100\% ranh giới mô hình frontend tải qua bộ đôi S3 và CloudFront CDN \\
Viết cẩu rằng kho S3 núp cẩn sau tấm khiên ALB & Lầm lọt mô hình thiết
định origin cõi CloudFront & Behavior mặc định của CloudFront buộc bão
dốc cầu trúng đích vào origin S3 thô trần \\
Xưng tụng nhầm cụm Redis chính là kho CSDL lưu chat vĩnh cửu & Hiểu sai
thảm lật cấu trúc luồng lưu trữ đàm thoại chat & Bàn DynamoDB chịu cọc
gánh việc ghi giữ tin nhắn trần; còn cõi Redis sắm quyền hạn bộ phát
sóng pub/sub \\
Miêu tả cơ quan outbox là phân phối đúng nhĩ chính xác 1 lần
(exactly-once) & Chưa am tường tập quy luật phát tin của SQS Standard &
Văn khải mô tả rõ tiêu chuẩn gửi ít nhất một lần (at-least-once) cõng
trói cùng thao thấu consumer có tính bất biến \\
Bồi cất kích bật công tắc worker trước lúc dàn hạ tầng AI sẵn sàng & Trụ
cọc thọ sinh SagMaker chưa mở thông & Ép giữ worker an bế trong nhãn
giam khóa tịt (disabled) đến mốc khoảnh khắc cụm AI cùng điểm endpoint
thọ thông bãi \\
\end{longtable}
\endgroup
}

\vspace{0.5em}\noindent\textbf{Kết luận}\par

Chương trình Kiến trúc này cống hiến mẫu mẫn khuôn bản tham chiếu nòng
cốt cho muôn chặng đường kiểm thử và thau dựng của Workshop kíp tiếp
bước. Mọi cương mục sau đó chuyên về ban bố triển khai, nghiệm thu test,
rà soát quan trắc, an nộ bảo mật, cước thọ ngân sách, xử trí lỗi lầm
cùng thao tác gõ dẹp tháo dỡ thảy đều bám đuổi tuyệt đối trung thành
theo nhĩ đàn sơ đồ tuần tự và hạ tầng đã tạc lên tại đây.
